\documentclass[11pt,a4paper]{article}
\usepackage[utf8]{inputenc}
\usepackage[T1]{fontenc}
\usepackage{tgpagella}
\usepackage[spanish]{babel}
\usepackage{geometry}
\usepackage{booktabs}
\usepackage{xcolor}
\usepackage{titlesec}
\usepackage{fancyhdr}
\usepackage[parfill]{parskip}
\usepackage[expansion=false]{microtype}
\usepackage{hyperref}
\usepackage{setspace}
\usepackage{needspace}
\usepackage{etoolbox}
\usepackage{graphicx}

\widowpenalty=10000
\clubpenalty=10000
\displaywidowpenalty=10000
\raggedbottom
\sloppy
\hbadness=10000
\hfuzz=2pt

\geometry{top=3.0cm,bottom=3.0cm,left=3.5cm,right=3.5cm}
\setlength{\parskip}{0.65em}
\setlength{\parindent}{0em}
\setlength{\headheight}{14pt}

\definecolor{azulai}{RGB}{30,80,140}
\definecolor{doradoai}{RGB}{140,90,0}
\definecolor{grisai}{RGB}{70,70,70}

\titleformat{\section}{\Large\bfseries\color{azulai}}{}{0em}{}[{\color{azulai}\titlerule[0.5pt]}]
\titlespacing*{\section}{0pt}{1.8em}{0.8em}

\titleformat{\subsection}{\large\bfseries\color{doradoai}}{}{0em}{}
\titlespacing*{\subsection}{0pt}{1.4em}{0.5em}

\titleformat{\subsubsection}{\normalsize\bfseries\color{grisai}}{}{0em}{}
\titlespacing*{\subsubsection}{0pt}{1.0em}{0.3em}

% Evita títulos huérfanos al final de página
\pretocmd{\section}{\Needspace{12\baselineskip}}{}{}
\pretocmd{\subsection}{\Needspace{9\baselineskip}}{}{}
\pretocmd{\subsubsection}{\Needspace{6\baselineskip}}{}{}

\pagestyle{fancy}
\fancyhf{}
\rhead{\textcolor{grisai}{\small Izas — Arquitectura Interna}}
\lhead{\textcolor{grisai}{\small Luna · Casa 4 · Casa 6}}
\cfoot{\textcolor{grisai}{\small\thepage}}
\renewcommand{\headrulewidth}{0.3pt}

\hypersetup{colorlinks=true,linkcolor=azulai,urlcolor=azulai}
\setstretch{1.45}
\tolerance=2500
\emergencystretch=8em

\begin{document}

% ── Portada ──────────────────────────────────────────────────────────────────
\begin{titlepage}
  \centering
  \vspace*{1.5cm}
  {\Huge\bfseries\color{azulai} Luna · Casa 4 · Casa 6}\\[0.5cm]
  {\large\color{grisai} Arquitectura Interna}\\[0.3cm]
  {\small\itshape\color{grisai} Un método para sostener cuerpo, energía y vida con coherencia}\\[2cm]
  {\huge\color{doradoai} Izas}\\[1.5cm]
  {\Large 16/04/1979 \quad 04:20}\\[0.3cm]
  {\Large Pamplona, España}\\[0.3cm]
  {\normalsize Lat: 42.8157° \quad Lon: -1.6522° \quad UTC+02:00}\\[0.3cm]
  {\normalsize Ascendente: Acuario 6°20' \quad
    MC: Escorpio 29°12'}\\[2cm]

  \begin{tabular}{lll}
    \textbf{Luna:} & Sagitario & Casa 10 \\
    \textbf{Casa 4:} & Tauro en cúspide & Regente: Venus \\
    \textbf{Casa 6:} & Cáncer en cúspide & Regente: Luna \\
  \end{tabular}\\[2cm]

  \vfill
  {\small Generado el 19/05/2026}
\end{titlepage}

\tableofcontents
\clearpage

% ── Datos de referencia ───────────────────────────────────────────────────────
\section{Datos de referencia}

\begin{center}
\begin{tabular}{llll}
  \toprule
  \textbf{Eje} & \textbf{Signo} & \textbf{Posición} & \textbf{Datos} \\
  \midrule
  Luna & Sagitario & Casa 10 & 8°24' \\
  Casa 4 (cúspide) & Tauro & --- & Regente: Venus \\
  Casa 6 (cúspide) & Cáncer & --- & Regente: Luna \\
  Ascendente & Acuario & --- & 6°20' \\
  Medio Cielo & Escorpio & --- & 29°12' \\
  \bottomrule
\end{tabular}
\end{center}

\vspace{0.4cm}
\textbf{Planetas en Casa 4:} ninguno\\
\textbf{Planetas en Casa 6:} Júpiter

\vspace{0.5cm}
\textbf{Aspectos de la Luna:}

\begin{center}
\begin{tabular}{lllll}
  \toprule
  & \textbf{Asp.} & \textbf{Planeta} & \textbf{Tipo} & \textbf{Orbe} \\
  \midrule
  Luna & □ & Saturno & Cuadratura & 0.8° \\
  Luna & ⚻ & Quirón & Quincuncio & 0.9° \\
  Luna & △ & Marte & Trígono & 1.4° \\
  \bottomrule
\end{tabular}
\end{center}

\vspace{0.8cm}
\begin{center}
\includegraphics[width=0.80\textwidth]{C:/Users/Izas/claude/astrología/pruebas/Izas_Luna_C4_C6_rueda.png}
\end{center}


\vspace{0.5cm}
\Needspace{5\baselineskip}

% ── Interpretación ────────────────────────────────────────────────────────────
\section{Interpretación — Arquitectura Interna}

\Needspace{8\baselineskip}
\begin{center}
{\small\itshape
No se trata de explicarte. Se trata de mostrar cómo funciona tu regulación interna,\\
dónde puede haber desregulación y qué necesitas para sostenerte.
}
\end{center}

\vspace{0.8cm}
\Needspace{5\baselineskip}

% ── 1. Luna ───────────────────────────────────────────────────────────────────
\subsection{1. Tu Luna — Función emocional}

\subsubsection*{Luna en Sagitario — Casa 10}

Necesitas sentir que lo que estás viviendo tiene dirección.

Cuando una emoción encuentra sentido o movimiento, puedes atravesarla mucho mejor. El problema aparece cuando sientes algo que no sabes hacia dónde va o que parece no conducir a ningún lugar.

Entonces suele aparecer inquietud, necesidad de escapar, cambiar de aire o salir de la situación de alguna manera. Tu cuerpo necesita espacio, movimiento y sensación de horizonte.

Cuando pasas demasiado tiempo sintiendo que no existe libertad o amplitud suficiente, la activación empieza a acumularse.

La regulación mejora cuando existe dirección, expansión y posibilidad de avanzar. En cambio, la sensación de límite, encierro o falta de sentido tiende a aumentar rápidamente la inquietud interna.

Lo que ocurre en tu vida profesional o pública te afecta emocionalmente mucho más de lo que aparentas. El reconocimiento puede darte una sensación profunda de estabilidad interna, mientras que la crítica o la exposición pueden impactarte muchísimo aunque no siempre lo muestres. Muchas veces la sensación de valor personal se conecta con cómo sientes que ocupas tu lugar frente al mundo. Cuando desaparece la dirección o el reconocimiento, la energía emocional puede caer muy rápido. Te ayuda sentir propósito, coherencia y reconocimiento genuino. En cambio, la exposición excesiva, la sensación de fracaso o la inseguridad pública suele afectar profundamente tu estabilidad emocional.

Puede costarte permitir ciertas emociones sin intentar controlarlas enseguida. La vulnerabilidad puede sentirse peligrosa, excesiva o difícil de mostrar abiertamente. Muchas veces una parte de ti intenta contener lo que siente antes incluso de haberlo reconocido por completo. No necesitas destruir tu estructura interna para sentir. Necesitas permitir que la emoción exista dentro de ella sin convertirla automáticamente en un problema que corregir.

Existe un ajuste constante entre lo que sientes y una sensibilidad emocional más profunda que se activa de formas diferentes según el momento. Lo que te ayuda a regularte en una etapa puede no servir igual en otra. No existe aquí una única forma fija de equilibrio emocional. Necesitas aprender a escucharte continuamente y adaptar el sostén a lo que realmente está ocurriendo dentro de ti en cada momento.

Tu emoción y tu capacidad de acción suelen trabajar juntas. Cuando algo te mueve por dentro, normalmente puedes transformarlo con bastante rapidez en movimiento, decisión o acción concreta. Eso ayuda a que la energía emocional no se quede completamente acumulada dentro del cuerpo. El movimiento y la acción consciente pueden ayudarte muchísimo a ordenar lo que sientes.

\vspace{0.8cm}
\Needspace{5\baselineskip}

% ── 2. Casa 4 ─────────────────────────────────────────────────────────────────
\subsection{2. Tu Casa 4 — Raíz interna}

\subsubsection*{Tauro en la cúspide — Regente: Venus}

Necesitas estabilidad para sentir seguridad interna. Los cambios bruscos en el hogar, en las condiciones de vida o en la sensación de seguridad te afectan mucho más de lo que suele verse desde fuera. Cuando algo importante cambia demasiado rápido, puede aparecer una sensación muy profunda de haber perdido suelo. Tiendes a buscar lo conocido, la repetición y cierta continuidad para volver a sentir calma. Por eso no suele resultar fácil relajarte cuando todo alrededor cambia constantemente. Te ayuda la calma, la estabilidad y sentir que existe un ritmo previsible en tu vida. En cambio, la incertidumbre, los cambios repentinos o sentir que pierdes tu base suele generar mucha inseguridad interna.

El regente de la Casa 4 es Venus, en Piscis y en Casa 1. La base interna tiende a organizarse desde la identidad y el cuerpo, con la cualidad receptiva y emocional de Piscis.

\vspace{0.8cm}
\Needspace{5\baselineskip}

% ── 3. Casa 6 ─────────────────────────────────────────────────────────────────
\subsection{3. Tu Casa 6 — Regulación cotidiana}

\subsubsection*{Cáncer en la cúspide — Regente: Luna}

Tu cuerpo y tu estado emocional están profundamente unidos. Cuando algo te afecta emocionalmente, el cuerpo suele notarlo enseguida. La digestión, el descanso, la energía o la sensación de agotamiento cambian rápidamente según cómo estén tus vínculos cercanos y el clima emocional que te rodea. Necesitas sentir cuidado y cierta seguridad emocional para sentirte bien físicamente de verdad. Los pequeños rituales cotidianos tienen muchísimo impacto sobre ti: la comida, el descanso, el hogar o la sensación de refugio. Cuando pasas demasiado tiempo sosteniendo tensión emocional, el cuerpo termina expresándolo de alguna manera. Te ayudan el cuidado cotidiano, la intimidad y sentir hogar alrededor. En cambio, la tensión relacional, la falta de descanso emocional o los ambientes afectivamente inestables suele desgastar rápidamente tu energía.

El regente de la Casa 6 es la Luna, en Sagitario y en Casa 10. La regulación cotidiana tiende a organizarse desde la vocación y el espacio público, con la cualidad activa y directa de Sagitario.

Te cuesta notar automáticamente cuándo algo ya es suficiente. Puedes trabajar más, hacer más o exigirte más de lo que realmente necesitas sin darte cuenta enseguida. No siempre ocurre por ambición consciente. Muchas veces simplemente no aparece el límite de forma natural. El problema es que el cuerpo sí tiene límite, aunque tardes en percibirlo. Necesitas construir conscientemente momentos de pausa, moderación y descanso. Cuando no existe esa sensación de medida, la expansión continúa hasta que el cuerpo empieza a pasar factura.

\vspace{0.8cm}
\Needspace{5\baselineskip}

% ── 4. Integración ────────────────────────────────────────────────────────────
\subsection{4. Integración — Cómo funcionan las tres capas juntas}

Tu Luna en Sagitario muestra cómo reaccionas emocionalmente cuando algo te afecta.

Es el primer movimiento: lo que aparece antes de que puedas ordenarlo, explicarlo o decidir qué hacer con ello.

Después, esa emoción busca un lugar donde apoyarse. Ahí entra tu Casa 4, con Tauro en la cúspide: la parte de ti que necesita sentir base, refugio y seguridad interna.

Si esa base está disponible, lo que sientes puede moverse con más facilidad.
Si no lo está, la emoción se queda más tiempo en estado de alerta.

Y lo que no encuentra apoyo interno termina llegando al cuerpo. Ahí entra tu Casa 6, con Cáncer en la cúspide: tus hábitos, tu energía diaria, tu descanso y la forma en que el cuerpo sostiene lo que vives.

Por eso, cuando algo se desordena emocionalmente, no se queda solo en lo emocional. Puede acabar apareciendo como cansancio, pérdida de ritmo, dificultad para descansar o sensación de no poder sostener el día con la misma energía.

En tu caso, las tres capas tienen elementos distintos: Fuego en la Luna, Tierra en la Casa 4 y Agua en la Casa 6.

Esto significa que tu emoción, tu base interna y tu cuerpo no siempre piden lo mismo.

Lo que te ayuda emocionalmente puede no ser lo que tu cuerpo necesita ese día. Lo que te da seguridad interna puede no coincidir con el ritmo que sostiene tu rutina.

Por eso es importante no buscar una única respuesta para todo. A veces una parte de ti necesita expresión, otra necesita refugio y otra necesita descanso, orden o movimiento.

Además, Júpiter en la Casa 6 influye en cómo tu cuerpo sostiene la vida cotidiana.

No se trata solo de hábitos o rutina. Se trata de cómo organizas tu energía, cómo respondes al cansancio y qué ocurre cuando el cuerpo empieza a pedir ajuste.

Esos planetas muestran qué tipo de fuerza, sensibilidad o tensión entra en tu regulación diaria.

Los aspectos de tensión con Saturno y Quirón hacen que el mundo emocional tenga menos margen en ciertos momentos.

Cuando algo te activa, puede haber menos tiempo entre sentir y reaccionar. La emoción puede llegar más intensa, más mezclada o más difícil de ordenar.

Esto no significa que sí o sí vayas a desregularte. Significa que necesitas reconocer antes las señales iniciales, porque cuando la activación ya ha subido mucho, cuesta más volver al centro.

Al mismo tiempo, los aspectos de apoyo con Marte ofrecen recursos reales.

Cuando consigues encontrar un primer punto de estabilidad, esos aspectos ayudan a sostenerlo.

No eliminan la tensión, pero pueden facilitar que vuelvas a organizarte sin quedarte completamente dentro de lo que se activó.

Lo más difícil no suele ser que una sola parte se altere.

Lo más difícil aparece cuando las tres capas se mueven al mismo tiempo:
la Luna en Sagitario, con su forma concreta de sentir;
la Casa 4 en Tauro, con su necesidad de base y seguridad interna;
y la Casa 6 en Cáncer, con el cuerpo intentando sostenerlo todo en la vida diaria.

Cuando esto ocurre, puedes notar que no sabes por dónde empezar: emocionalmente hay activación, internamente falta suelo y el cuerpo empieza a mostrar cansancio, tensión o pérdida de ritmo.

El trabajo no consiste en resolverlo todo a la vez.

Consiste en encontrar la primera capa disponible: una emoción que pueda expresarse, un entorno que pueda darte más seguridad, o una acción corporal sencilla que te ayude a volver poco a poco a la estabilidad.

\vspace{0.8cm}
\Needspace{5\baselineskip}

% ── 5. Sostén y desregulación ─────────────────────────────────────────────────
\subsection{5. Sostén y desregulación}

Tu Luna en Sagitario se desregula cuando lo que sientes no encuentra salida.

Puede ser porque no puedes expresarlo, porque no hay espacio para moverte, porque el entorno no lo recibe o porque intentas sostenerlo demasiado tiempo dentro.

No siempre hace falta que ocurra algo muy intenso.

A veces basta con acumular durante demasiado tiempo una emoción que no ha tenido lugar, palabra, movimiento o descanso.

Los aspectos de tensión con Saturno y Quirón hacen que ese margen sea más estrecho.

Cuando algo te activa, puede haber menos tiempo entre sentir y reaccionar.

La emoción puede llegar más fuerte, más mezclada o más difícil de ordenar. Por eso es importante reconocer las primeras señales antes de llegar al punto de desborde.

Tu Casa 4 con Tauro pierde estabilidad cuando dejas de sentir base interna.

No siempre tiene que haber un conflicto claro.

A veces basta con que cambie el clima emocional del entorno, con que una relación cercana se tense, con que no puedas descansar de verdad o con que el espacio íntimo deje de sentirse seguro.

Cuando eso ocurre, no solo cambia el ánimo. Cambia la sensación de tener suelo por dentro.

Tu Casa 6 con Cáncer muestra la desregulación en el cuerpo y en la vida cotidiana.

Puede aparecer como cansancio, pérdida de ritmo, tensión física, sueño menos reparador o dificultad para sostener hábitos que normalmente sí puedes sostener.

Muchas veces el cuerpo avisa antes de que puedas explicar qué está pasando.

Cuando el cuerpo empieza a perder continuidad, conviene mirar también qué emoción no ha encontrado salida y qué parte de tu base interna ha dejado de sentirse segura.

Los momentos más difíciles aparecen cuando las tres capas se mueven a la vez.

La Luna en Sagitario se activa.
La Casa 4 en Tauro deja de sentirse como base.
Y la Casa 6 en Cáncer empieza a mostrarlo en el cuerpo, el descanso o la rutina.

Ahí puedes sentir que no sabes por dónde empezar.

No se trata de resolverlo todo a la vez.
Se trata de encontrar la primera capa disponible: expresar algo, descansar, ordenar una pequeña rutina, moverte, pedir espacio o recuperar una mínima sensación de seguridad.

Los aspectos de apoyo con Marte funcionan como recursos cuando consigues recuperar un primer punto de estabilidad.

No eliminan la desregulación, pero pueden ayudarte a reorganizarte antes.

Cuando aparece una mínima base, esos aspectos facilitan que no te quedes dentro de lo que se activó.

\vspace{0.8cm}
\Needspace{5\baselineskip}

% ── 6. Orientación final ──────────────────────────────────────────────────────
\subsection{6. Orientación final}

\subsubsection*{Qué observar}

Empieza observando qué situaciones activan tu Luna en Sagitario.

No para corregirte.
No para controlar lo que sientes.

Solo para reconocer antes cómo empieza la activación.

¿Qué ocurre justo antes de sentirte que la situación te sobrepasa?
¿Qué cambia en el entorno, en los vínculos o en el ritmo diario?

Muchas veces la desregulación no empieza en el momento del desborde. Empieza bastante antes: cuando tu Casa 4 en Tauro deja de sentirse segura o cuando tu Casa 6 en Cáncer pierde continuidad y el cuerpo empieza a acumular tensión.

\subsubsection*{Qué puede desregularte}

En tu caso, la desregulación suele seguir una secuencia.

Primero, tu Luna en Sagitario busca algo que necesita para estabilizarse y no lo encuentra.

Después, tu Casa 4 deja de funcionar como base interna. Puede ser por tensión emocional, por cambios en el entorno, por exceso de carga o simplemente porque llevas demasiado tiempo sosteniendo más de lo que puedes absorber.

Y finalmente, la Casa 6 en Cáncer empieza a mostrarlo en el cuerpo:
menos energía,
menos descanso,
menos capacidad de sostener hábitos o continuidad.

Cuando las tres capas coinciden, puedes sentir que todo cuesta más de lo habitual aunque desde fuera no parezca estar ocurriendo nada tan grave.

\subsubsection*{Qué te estructura y sostiene}

Lo que te ayuda a regularte no tiene que ser enorme.

De hecho, muchas veces lo más importante son cosas pequeñas que se sostienen en el tiempo.

Una mínima rutina corporal.
Un espacio de descanso real.
Un entorno que se sienta más seguro.
Una forma de expresar lo que ocurre antes de acumularlo demasiado.

Cuando tu Casa 4 siente suelo, tu Luna necesita menos esfuerzo para estabilizarse.

Y cuando tu Casa 6 mantiene cierta continuidad —aunque sea pequeña— el cuerpo puede absorber mejor lo que emocionalmente estás atravesando.

No se trata de no sentir.
Se trata de que lo que sientes no tenga que terminar descargándose siempre en agotamiento, tensión o pérdida de equilibrio.

\subsubsection*{Práctica}

Tu regulación necesita vaciado emocional.

Dedica un momento del día a cerrar lo absorbido:
silencio,
respiración,
una ducha consciente,
música tranquila,
o simplemente estar sin estímulos.

No es una pausa superficial.
Es una forma de ayudarte a soltar lo que has ido acumulando.

Cuando no existe ese cierre, el cuerpo sigue cargando emociones incluso mientras intenta descansar.

\Needspace{8\baselineskip}
\vspace{0.6cm}

\begin{center}
{\small\itshape\color{grisai}
La astrología se usa aquí como lenguaje simbólico de observación, no como definición de la persona.\\
Lo que tienes delante es un mapa de posibilidades, no un destino fijo.
}
\end{center}

\end{document}
